\begin{proof}[Proof of \cref{obj:lem:regret-separation}]
Take the separation constant
\[
c_{\mathrm{sep}}=\frac{\min\{c_B,1\}}{4}>0,
\]
and work at all sufficiently large \(n\) for which \cref{obj:lem:witness-membership} applies to both signs, so that \(P_{n,+},P_{n,-}\in\mathcal P_{\alpha,\gamma}\) and their oracle rules are
\[
\pi^\star_{P_{n,+}}(x)=1
\qquad\text{and}\qquad
\pi^\star_{P_{n,-}}(x)=\mathbf 1\{x\notin B_n\}.
\]
% lean: witness_membership
Write \(h=h_n\), \(\tau_\pm=\tau_{P_{n,\pm}}\), and shrink the active block to
\[
E=[0,\ell],
\qquad
\ell=\min\{c_B,1\}\,h^\alpha .
\]
Since \(0<h\le1\) and \(\alpha\ge0\), we have \(0\le h^\alpha\le1\) and hence \(0\le\ell\le1\); as \(P_X\) is Lebesgue measure restricted to \([0,1]\), the covariate mass of \(E\) is exactly
\[
P_X(E)=\ell .
\]
% lean: restricted_volume_real_Icc_zero
Moreover \(\ell=\min\{c_B,1\}h^\alpha\le c_Bh^\alpha\), so
\[
E\subseteq B_n=\{x:0\le x\le c_Bh^\alpha\}.
\]

Fix any policy \(\pi\in\Pi\).  Let \(D_+\) and \(D_-\) be its disagreement sets under \(P_{n,+}\) and \(P_{n,-}\), and set
\[
B_+=D_+\cap\{x:h/2<|\tau_+(x)|\},
\qquad
B_-=D_-\cap\{x:h/2<|\tau_-(x)|\}.
\]
Under each of the two witness laws, which are well formed with bounded outcomes, the welfare identity of \cref{obj:thm:welfare-identity} writes the regret as \(\int|\tau|\mathbf 1_D\,dP_X\) with a nonnegative integrand exceeding \(h/2\) on the respective set \(B_\pm\); restricting the integral to that set therefore gives
\[
\frac h2\,P_X(B_+)\le R_{P_{n,+}}(\pi),
\qquad
\frac h2\,P_X(B_-)\le R_{P_{n,-}}(\pi).
\]
% lean: regret_disagreement_large_contrast_le

On \(E\subseteq B_n\), the contrasts are \(\tau_+(x)=h\) and \(\tau_-(x)=-h\), so both have magnitude \(h>h/2\); and the two oracle rules take opposite values there,
\[
\pi^\star_{P_{n,+}}(x)=1,
\qquad
\pi^\star_{P_{n,-}}(x)=0
\qquad (x\in E).
\]
Since \(\pi(x)\in\{0,1\}\), it differs from at least one of these two values at every \(x\in E\): if \(\pi(x)=1\) then \(x\in D_-\), and if \(\pi(x)=0\) then \(x\in D_+\).  Combined with the contrast bound this gives
\[
E\subseteq B_+\cup B_- .
\]
The two witness laws share the covariate marginal \(P_X\), so by monotonicity and subadditivity,
\[
\ell=P_X(E)\le P_X(B_+)+P_X(B_-).
\]
Multiplying by \(h/2>0\) and inserting the two regret bounds gives
\[
\frac h2\,\ell
\le \frac h2 P_X(B_+)+\frac h2 P_X(B_-)
\le R_{P_{n,+}}(\pi)+R_{P_{n,-}}(\pi).
\]
Since \(h^{1+\alpha}=h\cdot h^\alpha\) and \(\ell=\min\{c_B,1\}h^\alpha\), the left-hand side is
\[
\frac h2\,\ell=\frac{\min\{c_B,1\}}{2}\,h^{1+\alpha},
\qquad\text{so}\qquad
\frac{\min\{c_B,1\}}{2}h^{1+\alpha}
\le R_{P_{n,+}}(\pi)+R_{P_{n,-}}(\pi).
\]
Finally, a sum of two reals is at most twice their maximum,
\[
R_{P_{n,+}}(\pi)+R_{P_{n,-}}(\pi)
\le 2\max\{R_{P_{n,+}}(\pi),R_{P_{n,-}}(\pi)\},
\]
so dividing by \(2\),
\[
\max_{\sigma\in\{+,-\}}R_{P_{n,\sigma}}(\pi)
\ge \frac{\min\{c_B,1\}}{4}\,h_n^{1+\alpha}
= c_{\mathrm{sep}}\,h_n^{1+\alpha}.
\]
Since \(\pi\in\Pi\) was arbitrary, this holds for every policy in \(\Pi\), for all such sufficiently large \(n\). % lean: regret_separation
\end{proof}
